% Source audit for A.5, 15 November 2026.  Ready for root integration.
% No assertion that a corrected Sali\'e conjecture implies the actual arc count.
\paragraph{A refuted candidate for the conditional square-root argument.}
The literal Conjecture~1 of Bir\'o \cite{A5Biro2026} cannot serve as a
sound conjectural input: it contradicts a proved first-moment formula.
We give the specialization, including its smoothness and normalization
checks.  This does not refute a version in which the unary-theta
contribution is subtracted.

For odd positive $c$, put
\[
 T(m,n;c)=\sum_{x\in(\Z/c\Z)^\times}
 \left(\frac{x}{c}\right)e\left(\frac{m\bar x+nx}{c}\right),
 \qquad e(u)=e^{2\pi i u},
 \qquad
 \epsilon_c=\begin{cases}1,&c\equiv1\pmod4,\\
                         i,&c\equiv3\pmod4.
            \end{cases}
\]
These are Bir\'o's conventions.  Louvel's $K_2(m,n;c)$ writes the two
arguments in the reverse order; replacing $x$ by $\bar x$ shows that
the two definitions agree, since
$\bigl(\frac{\bar x}{c}\bigr)=\bigl(\frac{x}{c}\bigr)$.

Choose a fixed nonnegative, nonzero function
$g\in C_c^\infty(\R)$ supported in $[1,2]$ and, for $C\geq1$, define
\[
 f_C(x)=\sqrt{C/x}\,g(x/C).
\]
This is smooth and supported in $[C,2C]$, and for every fixed
$j\geq0$ it satisfies
$x^j f_C^{(j)}(x)=O_{g,j}(1)$.
In Bir\'o's conjecture take
\[
 L=4,\quad K=1,\quad m=4,\quad n=1,
 \quad B=1,\quad\alpha=0,
 \quad r=1.
\]
All the hypotheses hold: $2\mid L$,
$\gcd(L,Kr)=1$, and $LK(1+B)=8\ll_\delta C^\delta$
for any $\delta>0$ once $C$ grows.  The stated derivative bounds hold
for that same $\delta$.  For each odd $c$, multiplication by $4$ is a
bijection of $(\Z/c\Z)^\times$.  Substitution $x=4y$ gives
\[
 T(4,\bar4;c)
 =\sum_y\left(\frac{4y}{c}\right)
 e\left(\frac{\bar y+y}{c}\right)
 =T(1,1;c).
\]
The choice $n=1$ also removes any possible notational ambiguity about
the placement of the overline in Bir\'o's second argument.
Thus his asserted estimate would give
\[
 \sum_{c\equiv1\pmod4}
 g(c/C)\frac{T(1,1;c)}{\sqrt c}
 =O_{g,\varepsilon}(C^{1/2+\varepsilon}).
\]

Theorem~1 of Louvel's published paper
\cite[p.~525]{A5Louvel2012} (Theorem~1.1 in its arXiv version)
has the following applicable specialization:
\begin{equation}\label{A5audit:louvel-specialization}
 A(U):=\sum_{\substack{0<c\leq U\\c\ \text{ odd}}}
 \frac{T(1,1;c)\overline{\epsilon_c}}{\sqrt c}
 =\frac8{\pi^2}U+O_\eta(U^{3/4+\eta})
 \qquad(\eta>0).
\end{equation}
For completeness, in his general theorem we take the odd, coprime
integers $D=f=1$, the even primitive character of conductor~$1$,
and positive $m=n=1$.  In his formula for the main coefficient take
$t=s=m'=n'=D_t=D_0=D_1=1$.
Then $t$ is squarefree and $1$ modulo~$4$, $s$ is supported by $t$,
$s^2t^3\mid D$, and all coprimality and equal-gcd conditions hold.
The Jacobi symbol with denominator~$1$ equals~$1$ and all prime
products are empty, so the coefficient is exactly $8/\pi^2$.
The conjugate on $\epsilon_c$ has been checked in both the author
preprint and the published formula; omitting it would change this
normalization.

The same theorem, with $m=n=-1$ and the other data unchanged, gives
\[
 A_-(U):=\sum_{\substack{0<c\leq U\\c\ \text{ odd}}}
 \frac{T(-1,-1;c)\overline{\epsilon_c}}{\sqrt c}
 =O_\eta(U^{3/4+\eta}),
\]
because its first case applies whenever either argument is negative.
The exact substitution $x=-y$ in the defining sum gives
\[
 T(-1,-1;c)=\left(\frac{-1}{c}\right)T(1,1;c).
\]
Consequently the single residue class $1$ modulo $4$ has the asymptotic
\begin{equation}\label{A5audit:single-class}
 A_1(U):=\sum_{\substack{0<c\leq U\\c\equiv1\pmod4}}
 \frac{T(1,1;c)}{\sqrt c}
 =\frac{A(U)+A_-(U)}2
 =\frac4{\pi^2}U+O_\eta(U^{3/4+\eta}).
\end{equation}

Abel summation applied to \eqref{A5audit:single-class} gives
\begin{align*}
 \sum_{c\equiv1\pmod4} g(c/C)
       \frac{T(1,1;c)}{\sqrt c}
 &=-\frac1C\int_C^{2C} A_1(u)g'(u/C)\,du\\
 &=\frac4{\pi^2}C\int_1^2g(v)\,dv
       +O_{g,\eta}(C^{3/4+\eta}).
\end{align*}
In the exact normalization of Bir\'o's left side, this states
\[
 \frac1C\sum_{c\equiv1\pmod4}f_C(c)T(4,\bar4;c)
 =\frac4{\pi^2}\sqrt C\int_1^2g(v)\,dv
 +O_{g,\eta}(C^{1/4+\eta}).
\]
The main coefficient is positive.  Taking, for example,
$\eta=1/16$ and $\varepsilon=1/8$ contradicts his asserted
$O(C^\varepsilon)$ bound as $C\to\infty$.
This is a single completely specified parameter family violating the
printed conjecture.

The program \path{code/B_A5_salie_audit.py} checks the change of
variables and the identity for $(-1,-1)$ exactly as equalities of integer coefficient vectors in
$\Z[X]/(X^c-1)$ for every odd $3\leq c\leq1001$.
It also checks the identity
\[
 T(1,1;c)=
 \left(\sum_{y\bmod c}e(y^2/c)\right)
 \sum_{v^2\equiv4\pmod c}e(v/c)
\]
exactly modulo the cyclotomic polynomial $\Phi_c$ for odd
$3\leq c\leq199$.
The accompanying first-moment evaluations use 192-bit Arb arithmetic
and the explicit bump
\[
 g(v)=\begin{cases}
 \exp\!\left(4-\dfrac1{(v-1)(2-v)}\right),&1<v<2,\\
 0,&\text{otherwise}.
 \end{cases}
\]
They are finite normalization checks, not the proof of the asymptotic
contradiction.  The latter is the preceding application of Louvel's
theorem.  On $[5/4,7/4]$ this bump is at least $e^{-4/3}$, so its
positive integral does not depend on a numerical quadrature.

For the literature status of this correction, on 15 November 2026 we
searched the exact paper identifier, its title together with
``correction'' and ``counterexample'', and Bir\'o--Louvel--Sali\'e
combinations.  The primary arXiv record still listed only version~1
of 13 April 2026.  We did not locate an earlier published correction
of this printed conjecture.  The mathematical input is Louvel's known
theorem, and the present deduction is an independently checked
specialization of it; we make no claim of priority.

\paragraph{The corrected candidate and the remaining reduction.}
Conjecture~1.4 of Sun \cite{A5Sun2025} explicitly subtracts
\[
 2Y^{1/2}\sum_{r_j=i/4}\tau_j(m,n)
\]
from the corresponding sum of half-integral-weight Kloosterman
sums divided by their moduli.  Its hypotheses include a weight
$k=\pm1/2$ multiplier on $\Gamma_0(N)$ and absence of Laplace
eigenvalues in $(3/16,1/4)$ for that multiplier space.
It gives an untwisted, uniformly parameterized modulus-sum estimate;
it does not state a bound for the angular-period coefficients of
our compact quaternionic surface.
We have not established a reduction from that conjecture to the
actual sharp-arc count.  In a direct ternary delta-method approach
one must specify the multiplier and its spectral condition, handle
the local factors at $2$ and $3$, identify and subtract every theta
term, and prove uniform bounds for the resulting modulus weights
as the physical angular and radial transition widths shrink to
$X^{-1/2}$.  Partial summation with an uncontrolled total-variation
factor does not prove the required exponent.
